\section{TaxoNet Architecture}

\label{sec:architecture}
% ══════════════════════════════════════════════════════════════════════════════

\subsection{Overview}

TaxoNet consists of four components: (1) a vision transformer feature
extractor, (2) a taxonomic prototype constructor, (3) a rank-attention
fusion module, and (4) optional geographic/phenological prior
integration.

\subsection{Feature Extractor}

We use a Vision Transformer (ViT-L/14) pre-trained through our
domain-adaptive pipeline (Section~\ref{sec:pretraining}). Input images
are resized to $518 \times 518$ pixels and split into $14 \times 14$
patches. The transformer produces a sequence of patch embeddings, from
which we extract both the \texttt{[CLS]} token embedding $\mathbf{h}
\in \R^{1024}$ and the full patch embedding matrix $\mathbf{H} \in
\R^{1369 \times 1024}$.

For fine-grained discrimination, we apply a learned spatial attention
module that identifies discriminative image regions:
\begin{equation}
  \mathbf{a} = \text{softmax}(\mathbf{W}_a \mathbf{H}^\top) \in
  \R^{M \times 1369}
\end{equation}
where $M = 4$ attention heads focus on different discriminative parts.
The attended feature is:
\begin{equation}
  \mathbf{f} = [\mathbf{h} \| \text{vec}(\mathbf{a} \mathbf{H})]
  \in \R^{1024 + 4 \times 1024} = \R^{5120}
\end{equation}
projected to a 512-dimensional embedding via a linear layer.

\subsection{Taxonomic Prototypical Network}

Given the support set $\Support$ for an $N$-way $K$-shot episode, we
construct prototypes at three taxonomic ranks: family ($r_F$), genus
($r_G$), and species ($r_S$).

\paragraph{Species prototype.} The standard prototypical network
prototype:
\begin{equation}
  \mathbf{c}_{r_S}^{(s)} = \frac{1}{K} \sum_{(x_i, y_i) \in
  \Support : y_i = s} \phi(x_i)
  \label{eq:species_proto}
\end{equation}
where $\phi(\cdot)$ is the feature extractor.

\paragraph{Genus prototype.} Aggregating all species in the same genus:
\begin{equation}
  \mathbf{c}_{r_G}^{(g)} = \frac{1}{|\{s : t_{r_G}(s) = g\}|}
  \sum_{s : t_{r_G}(s) = g} \mathbf{c}_{r_S}^{(s)}
\end{equation}

\paragraph{Family prototype.} Similarly aggregating across genera in
the same family:
\begin{equation}
  \mathbf{c}_{r_F}^{(f)} = \frac{1}{|\{g : t_{r_F}(g) = f\}|}
  \sum_{g : t_{r_F}(g) = f} \mathbf{c}_{r_G}^{(g)}
\end{equation}

\subsection{Rank-Attention Fusion}

For a query embedding $\mathbf{q} = \phi(x)$, we compute distances to
each prototype at each rank. The rank-attention module learns to weight
these distances adaptively:

\begin{equation}
  d_r(x, s) = \|\mathbf{q} - \mathbf{c}_r^{(t_r(s))}\|^2
\end{equation}

The attention weights depend on the query itself:
\begin{equation}
  \boldsymbol{\alpha}(\mathbf{q}) = \text{softmax}(\mathbf{W}_r
  \mathbf{q}) \in \R^3
\end{equation}

The fused distance is:
\begin{equation}
  d_{\text{fused}}(x, s) = \sum_{r \in \{r_F, r_G, r_S\}}
  \alpha_r(\mathbf{q}) \cdot d_r(x, s)
\end{equation}

The predicted class probability is:
\begin{equation}
  p(y = s \mid x) = \frac{\exp(-d_{\text{fused}}(x, s))}
  {\sum_{s'} \exp(-d_{\text{fused}}(x, s'))}
  \label{eq:prediction}
\end{equation}

\subsection{Geographic and Phenological Priors}

When location and date metadata are available, we compute a prior
probability $p_{\text{geo}}(s \mid \text{lat}, \text{lon}, d)$ from
species range maps and phenological data. This prior is combined with
the visual prediction via:
\begin{equation}
  p_{\text{final}}(s \mid x, \text{lat}, \text{lon}, d) \propto
  p(y = s \mid x)^{1-\beta} \cdot
  p_{\text{geo}}(s \mid \text{lat}, \text{lon}, d)^{\beta}
\end{equation}
where $\beta = 0.3$ balances visual and geographic evidence. This
effectively re-ranks predictions by geographic plausibility.


% ══════════════════════════════════════════════════════════════════════════════
